\chapter{De ITM-unit: massa- en energiebalans}
\label{h:itm}

\section{Gegevens uit de datasheet}
De ITM Neptune V is voor dit verslag een ``zwarte doos'' met de volgende gegevens per unit
(datasheet, nominaal bedrijf):
\begin{center}
\begin{tabular}{p{6.4cm}rp{4.6cm}}
\toprule
grootheid & waarde & stroom\\
\midrule
waterstofproductie $\md_{\HH}$ & \VAR{n(I.m_H2,0)}\kgh & S4, 25\degC, 30\,barg\\
zuurstofstroom inclusief water $\md_{\OO,\mathrm{mix}}$ & \VAR{n(I.m_O2mix,0)}\kgh & S13, 70\degC, 0,1\,barg\\
voedingswater totaal & \VAR{n(I.m_voedwater,0)}\kgh & S1 + S16\\
afvalwater & \VAR{n(I.m_afvalwater,0)}\kgh & S3\\
specifiek energieverbruik BOL & \VAR{n(I.SEC_BOL,1)}\,kWh/kg\,\HH & \\
restwarmte naar het koelcircuit (EOL) & \VAR{n(I.Q_ITM_EOL,0)}\kW & circuit A/B\\
koelcircuit & \VAR{n(I.m_kc_ds,0)}\kgh, \VAR{n(I.T_kc_retour,0)} $\rightarrow$ \VAR{n(I.T_kc_uit,1)}\degC & 62/38 water/CoolFlow DTX\\
massafractie water in S13 $w_{\HHO}$ & \VAR{n(I.w_H2O,2)} & aanname, zie §\ref{s:stoich}\\
\bottomrule
\end{tabular}
\end{center}
\begin{waarom}[BOL en EOL]
BOL (\emph{begin of life}) is de nieuwe stack; EOL (\emph{end of life}) is de verouderde stack vlak voor
vervanging. Een oude stack heeft een hogere celspanning bij dezelfde stroom, dus meer elektrisch vermogen voor
dezelfde hoeveelheid \HH. Het verschil wordt warmte. Voor de warmteterugwinning rekenen we met EOL: dan is er de
meeste warmte, en de apparatuur (E-201, TW, warmtepompen) moet die op het slechtste moment kunnen verwerken.
\end{waarom}

\section{Massabalans van de unit}
\kop{Gegeven}
In: voedingswater \VAR{n(I.m_voedwater,0)}\kgh. Uit: \HH\ \VAR{n(I.m_H2,0)}\kgh\ + \OO-mengsel
\VAR{n(I.m_O2mix,0)}\kgh\ + afvalwater \VAR{n(I.m_afvalwater,0)}\kgh.
\kop{Invullen}
\[
\sum \md_\mathrm{uit} = \VAR{n(I.m_H2,0)} + \VAR{n(I.m_O2mix,0)} + \VAR{n(I.m_afvalwater,0)} = \VAR{n(R.m_massa_uit,0)}\kgh
\quad\text{tegen}\quad \sum \md_\mathrm{in} = \VAR{n(R.m_massa_in,0)}\kgh .
\]
\begin{resultaat}
De datasheet sluit op \VAR{n(R.m_massa_in - R.m_massa_uit,0)}\kgh\ na
(\VAR{pct((R.m_massa_in - R.m_massa_uit)/R.m_massa_in, 2)}). Dat is een afrondings- of meetverschil in de datasheet
en niet iets wat het model veroorzaakt. Dit verschil komt terug in de controle van hoofdstuk~\ref{h:balans}.
\end{resultaat}

\section{Controle met de reactievergelijking}
\label{s:stoich}
\begin{methode}[– Cengel H15, massabalans van een reactie]
Bij elektrolyse geldt $\HHO \rightarrow \HH + \tfrac12\,\OO$. Per kmol \HH\ ontstaat dus een halve kmol \OO.
\end{methode}
\kop{Invullen}
\[
N_{\HH} = \frac{\VAR{n(I.m_H2,0)}}{\VAR{n(R.M_H2,3)}} = \VAR{n(R.N_H2_st,2)}\,\mathrm{kmol/h},\qquad
\md_{\OO,\mathrm{stoich}} = \tfrac12\cdot\VAR{n(R.N_H2_st,2)}\cdot\VAR{n(R.M_O2,3)} = \VAR{n(R.m_O2_st,1)}\kgh .
\]
Het water dat de reactie verbruikt is $\VAR{n(R.N_H2_st,2)}\cdot\VAR{n(R.M_H2O,3)} = \VAR{n(R.m_w_reactie,1)}\kgh$.
\begin{letop}[aandachtspunt voor de verdediging]
Het model neemt $w_{\HHO} = \VAR{n(I.w_H2O,2)}$. Dan bevat S13 $\VAR{n(R.m_O2,2)}\kgh$ \OO. De reactie levert bij
\VAR{n(I.m_H2,0)}\kgh\ \HH\ echter maar \VAR{n(R.m_O2_st,1)}\kgh\ \OO. Als de datasheetwaarde van
\VAR{n(I.m_O2mix,0)}\kgh\ klopt, bevat S13 dus $\VAR{n(I.m_O2mix,0)} - \VAR{n(R.m_O2_st,1)} = \VAR{n(R.m_H2O_st,1)}\kgh$
water: $w_{\HHO} \approx \VAR{n(R.w_st,3)}$ in plaats van \VAR{n(I.w_H2O,2)}.
\textbf{Gevolg:} met de stoichiometrische waarde komt er meer water mee naar E-101. Dan wordt
$\Qd_{\text{E-101}}$ groter en $\md_\mathrm{cond}$ groter; het dauwpunt van S13 stijgt ook. In bijlage~\ref{b:vragen} staat het effect
doorgerekend (vraag ``wat als $w_{\HHO} = 0{,}228$?''). De waarde 0,15 is een invoeraanname van het model. Noem dit
zelf als iemand vraagt of de \OO-stroom klopt.
\end{letop}

\section{Energiebalans van de unit}
\begin{methode}[– Cengel H5 + H15]
Neem de ITM-unit als controlevolume in stationaire toestand. Er gaat elektrische energie in ($P_\mathrm{AC}$).
Er gaat uit:
\begin{enumerate}
\item chemische energie in de waterstof: $\md_{\HH}\cdot\mathrm{HHV}$;
\item warmte naar het koelcircuit: $\Qd_\mathrm{ITM}$;
\item warmte die met de warme, natte \OO-stroom meegaat en in E-101 wordt afgevoerd: $-\Qd_{\text{E-101}}$.
\end{enumerate}
Dus $P_\mathrm{AC} = \md_{\HH}\,\mathrm{HHV}/3600 + \Qd_\mathrm{ITM} - \Qd_{\text{E-101}}$.
\end{methode}
\begin{aanname}
De voelbare warmte van \HH\ (25\degC), van het afvalwater en het verschil in temperatuur van het voedingswater worden
verwaarloosd. Dat zijn termen van enkele kW tegen ruim 5000\kW. Het referentiepunt is vloeibaar water bij
omgevingstemperatuur, en daarom is de HHV de juiste verbrandingswaarde.
\end{aanname}

\kop{Stap 1: chemische energie in de waterstof}
\[
\Qd_{\HH} = \frac{\md_{\HH}}{3600}\cdot\mathrm{HHV} = \frac{\VAR{n(I.m_H2,0)}}{3600}\cdot 141\,800 = \VAR{n(R.E_H2,1)}\kW .
\]
\begin{waarom}[HHV en niet LHV]
Het water gaat \emph{vloeibaar} de stack in. Dan moet de elektrolyse ook de verdampingswarmte leveren, en daarom
hoort de HHV (product vloeibaar water) bij elektrolyse. Met de LHV zou de \HH\ maar
$\VAR{n(R.E_H2_LHV,1)}\kW$ ``bevatten'' en zou er schijnbaar \VAR{n(R.E_H2 - R.E_H2_LHV,0)}\kW\ meer
warmte zijn, die er in werkelijkheid niet is.
\end{waarom}

\kop{Stap 2: elektrisch vermogen BOL (uit de datasheet)}
\[
P_\mathrm{AC,BOL} = \mathrm{SEC}\cdot \md_{\HH} = \VAR{n(I.SEC_BOL,1)}\,\mathrm{kWh/kg}\cdot\VAR{n(I.m_H2,0)}\kgh = \VAR{n(R.P_AC_BOL,0)}\kW .
\]
De bijbehorende restwarmte (BOL) volgt uit de balans:
$\Qd_\mathrm{ITM,BOL} = \VAR{n(R.P_AC_BOL,0)} - \VAR{n(R.E_H2,1)} + (\VAR{n(R.Q_E101,1)}) = \VAR{n(R.Q_ITM_BOL,1)}\kW$.

\kop{Stap 3: elektrisch vermogen EOL (uit de EOL-restwarmte)}
\[
P_\mathrm{AC,EOL} = \VAR{n(R.E_H2,1)} + \VAR{n(I.Q_ITM_EOL,0)} - (\VAR{n(R.Q_E101,1)}) = \VAR{n(R.P_AC_EOL,1)}\kW
\]
De waarde van $\Qd_{\text{E-101}}$ wordt afgeleid in hoofdstuk~\ref{h:o2}.
Dat geeft een specifiek verbruik van $\VAR{n(R.P_AC_EOL,1)}/\VAR{n(I.m_H2,0)} = \VAR{n(R.SEC_EOL,2)}$\,kWh/kg en een
rendement op HHV-basis van $\eta_\mathrm{HHV} = \VAR{n(R.E_H2,1)}/\VAR{n(R.P_AC_EOL,1)} = \VAR{pct(R.eta_HHV_EOL)}$
(BOL: \VAR{pct(R.eta_HHV_BOL)}).

\begin{resultaat}
Per unit is $P_\mathrm{AC,EOL} = \VAR{n(R.P_AC_EOL,0)}\kW$. Voor 8 units is dat
$\VAR{n(R.P_AC_tot,0)}\kW \approx \VAR{n(R.P_AC_tot/1000,1)}$\,MW. Van elke unit gaat \VAR{n(I.Q_ITM_EOL,0)}\kW\
restwarmte naar het koelcircuit, \VAR{pct(I.Q_ITM_EOL/R.P_AC_EOL)} van de elektrische input. Die warmte is de
belangrijkste bron voor het warmtenet.
\end{resultaat}
